% META: {"id":"1SPE-DERGLOBAL-EX-030","chapitre":"1SPE-DERIVATION-GLOBAL","type_objet":"exercice","capacites":["1SPE-DERIVATION-GLOBAL-C3"],"capacites_codes":["C3"],"methodes":["M3"],"parcours":3,"competences":["calculer","raisonner","communiquer"],"duree_min":20,"mode_creation":"ex_nihilo","sources_inspiration":[],"fichier_tex":"chapitres/1SPE-DERIVATION-GLOBAL/exercices/1SPE-DERGLOBAL-EX-030.tex","corrige_tex":"chapitres/1SPE-DERIVATION-GLOBAL/corriges/1SPE-DERGLOBAL-CO-030.tex","status":"approved","origin":"generated","status_history":["generated","audited","accepted","approved"],"release_acceptance":"RELEASE_OWNER_BATCH_ACCEPTANCE","acceptance_closure_digest":"sha256:9b3ccf9a81c5520fbb7e03b7b2d3e7bf2057a02834b6d908bf3be5f49f3b553f","acceptance_receipt_digest":"sha256:4fad86f0282e0fa4e0419cca1ad6379ae7af5a280c04a4a90880f90a1c044242"}
% BEGIN-VERIFY
% from sympy import symbols, diff, solve, Rational, sqrt
% x = symbols('x', real=True)
% f = 2*x**3 - 3*x**2 - 12*x + 4
% fp = diff(f, x)
% assert fp == 6*x**2 - 6*x - 12
% racines = solve(fp, x)
% assert set(racines) == {-1, 2}
% assert fp.subs(x, -2) > 0
% assert fp.subs(x, 0) < 0
% assert fp.subs(x, 3) > 0
% assert f.subs(x, -1) == 11
% assert f.subs(x, 2) == -16
% END-VERIFY
\begin{exercice}{1SPE-DERGLOBAL-EX-030}{3}{20}
Soit $f(x) = 2x^3 - 3x^2 - 12x + 4$ définie sur $\mathbb{R}$.

\begin{enumerate}
\item Calculer $f'(x)$ et montrer que $f'(x) = 6(x+1)(x-2)$.
\item Dresser le tableau de signes de $f'(x)$ puis le tableau de variations complet de $f$ sur $\mathbb{R}$.
\item \textbf{Preuve de monotonie.} Démontrer que $f$ est strictement décroissante sur $[-1\,;\,2]$ en utilisant la définition : pour tous $x_1, x_2 \in [-1\,;\,2]$ avec $x_1 < x_2$, montrer que $f(x_1) > f(x_2)$.

\textit{Indication : on pourra utiliser le théorème liant signe de la dérivée et sens de variation.}
\item Résoudre graphiquement (ou par calcul) l'inéquation $f'(x) \geqslant 0$.
\end{enumerate}
\end{exercice}
